\begin{table}[H]
\centering
\caption{Summary Statistics: Quarterly Employment Rates (\%), Poland 2010--2024}
\label{tab:summary}
\begin{threeparttable}
\begin{tabular}{lccccc}
\toprule
 & Quarters & Mean & SD & Min & Max \\
\midrule
\multicolumn{6}{l}{\textit{Panel A: Pre-reform (2010-Q1 to 2017-Q3)}} \\
Women 60--64 (treated) & 31 & 15.9 & 3.3 & 12.3 & 23.6 \\
Women 55--59 (control) & 31 & 47.0 & 8.0 & 30.2 & 59.3 \\
Men 60--64 (control) & 31 & 36.3 & 6.2 & 25.2 & 47.3 \\
Men 55--59 (control) & 31 & 64.8 & 3.5 & 57.6 & 71.3 \\
\midrule
\multicolumn{6}{l}{\textit{Panel B: Post-reform (2017-Q4 to 2024-Q4)}} \\
Women 60--64 (treated) & 29 & 22.8 & 2.3 & 19.1 & 27.8 \\
Women 55--59 (control) & 29 & 66.5 & 5.0 & 58.1 & 74.4 \\
Men 60--64 (control) & 29 & 56.6 & 5.7 & 46.1 & 64.1 \\
Men 55--59 (control) & 29 & 75.9 & 3.2 & 70.2 & 80.6 \\
\bottomrule
\end{tabular}
\begin{tablenotes}[flushleft]
\small
\item \textit{Notes:} Quarterly employment rates from Eurostat Labour Force Survey (\texttt{lfsq\_ergan}), Poland, all citizens. Pre-reform: 2010-Q1 to 2017-Q3 (31 quarters). Post-reform: 2017-Q4 to 2024-Q4 (29 quarters). The treated group (women 60--64) became newly eligible for statutory pension at age 60 after October 1, 2017. Men 60--64 remained below the retirement threshold (65) under both regimes.
\end{tablenotes}
\end{threeparttable}
\end{table}
